\chapter{主动热信息探针}
\label{ch:ati}

被动传感器在第 \ref{ch:est} 章已经用尽之后，剩余的信息只能来自有意改变边界条件。主动热信息（Active Thermal Intelligence, ATI）把蒸汽、热风、加热功率和排汽阀定义成一次短时实验，而不是烹饪程序的又一个阶段。

\section{探针与加工的区别}

\begin{definition}[热信息探针]
探针是一段有限时长 $\tau\le\tau_{\max}$ 的执行器轨迹 $a\in\cA$，满足：
\begin{enumerate}
\item 设计目标是增大 \eqref{eq:fisher} 的 $\log\det$，或等价地减小决策量的后验熵；
\item 对 $\alpha_b$、质量损失和 $\alpha_m$ 的预测损伤低于单独的预算 $\eta$；
\item 结束后执行器回到探针前的名义轨迹，除非安全层已经改写该轨迹。
\end{enumerate}
不满足第 1 条的功率变化是加工动作，归第 \ref{ch:safe} 章调度。
\end{definition}

同一阀门因此有两个目标函数。湿度调节最小化湿度误差。信息动作最大化参数信息。日志记录该时刻哪一个函数产生了阀门指令，执行器本身不区分来源。

\section{信息目标}

高斯近似下，决策状态的微分熵由协方差决定，
\begin{equation}
H(\vect{z}\mid\vect{y}_{0:k})=\frac12\log\big((2\pi\mathrm{e})^{n_z}\det\vect{\Sigma}_{z,k}\big).
\end{equation}
对候选探针 $a$，用当前后验抽样向前仿真到 $k+\tau$，得到预测观测分布与更新后的 $\vect{\Sigma}_{z,k+\tau}(a)$。信息增益
\begin{equation}
\label{eq:ati}
G(a)=\frac12\log\det\vect{\Sigma}_{z,k}
-\frac12\log\det\vect{\Sigma}_{z,k+\tau}(a).
\end{equation}
实现时用 $N$ 个集合成员做无迹变换或蒙特卡洛，不必假设全程高斯。与 Fisher 信息的关系是：在无偏估计的 Cramér–Rao 界附近，
\begin{equation}
G(a)\approx\frac12\log\det\vect{\mathcal{I}}(\vect{\theta};a)
-\frac12\log\det\vect{\mathcal{I}}(\vect{\theta};\emptyset),
\end{equation}
其中 $\emptyset$ 表示维持当前输入。在线计算使用 \eqref{eq:ati} 的集合形式。Fisher 形式在线性化工作点上平行计算，二者排序不一致时以集合形式为准，并记录该次失配。

代价取预测损伤与能量：
\begin{equation}
\label{eq:cost}
C(a)=w_E E(a)+w_m\Delta m(a)+w_b\Delta\alpha_b(a)+w_\tau\tau(a).
\end{equation}
最优探针
\begin{equation}
\label{eq:astar}
a^\star=\argmax_{a\in\cA}\Big(G(a)-\lambda C(a)\Big),
\end{equation}
仅当 $G(a^\star)\ge G_{\min}$ 且第 \ref{ch:est} 章的失效条件成立时执行。$\lambda$ 与 $G_{\min}$ 按食材类别标定：薄片蔬菜的 $G_{\min}$ 高（本来就容易观察，少探针），大块肉的 $G_{\min}$ 低。

\section{探针库，而不是连续最优控制}

家电不能在每个控制周期求解任意波形。探针库是有限集合，在线只做打分。

\begin{table}[H]
\centering
\caption{第一版探针库}
\begin{tabularx}{\textwidth}{lXll}
\toprule
编号 & 轨迹 & 主要敏感参数 & 时长 \\
\midrule
P1 & 蒸汽开、风机低、加热保持 & $K_{\mathrm{cd}}$、表面 $C_1$ & 4--8 s \\
P2 & 蒸汽关、排汽开、风机高 & $R_{\mathrm{ext}}$、$k_m$ & 6--10 s \\
P3 & 加热阶跃 $+15\%$、风机关 & 辐射份额、表面发射率路径 & 8--12 s \\
P4 & P1 后立即 P2 & 把冷凝跃迁与对流热阻分开 & $\le 20$ s \\
P0 & 维持现状 & 作为 $G$ 的零点 & 0 \\
\bottomrule
\end{tabularx}
\end{table}

P4 是默认的信息探针。单独的 P1 会把冷凝潜热和内部导热缠在一起；紧接着的排汽—高风只改变外部热阻与蒸发方向，两次响应的差对 $R_{\mathrm{int}}$ 敏感。这是本设计相对于“偶尔喷一点蒸汽看湿度”的具体机制。

P4 的机制可以收成一条可检查的规则。先使界面相变率变号，再使外部对流热阻下降且界面转为蒸发。两段采样之差更新 $R_{\mathrm{int}}$ 与 $C_2$。后验方差未下降则该次脉冲无效，停止时刻保持更新前的值。

\section{损伤预算如何算}

探针不是“短所以无害”。用 \eqref{eq:kinetics} 在表面节点上积分 $\Delta\alpha_b$，用 \eqref{eq:F} 在冷点上确认 $\alpha_m$ 的增量。若冷点在探针期间的 $\Delta F$ 可以忽略（秒级蒸汽对大块食品的冷点通常如此），预算主要约束表面褐变与冷凝水导致的表面浸湿。浸湿通过 $\varepsilon_{l,s}$ 的上升计入，后续烤制若需要干表面，则降低 P1 的允许时长。

存在硬否决，不进入 \eqref{eq:astar}：

\begin{itemize}
\item 水位低于干烧阈值；
\item 门打开；
\item 腔压或门封温度已接近安全边界；
\item 表面温度已高于褐变失控区，再做 P3。
\end{itemize}

\section{在线算法}

\begin{enumerate}
\item 每个估计周期计算信息赤字 $\Delta\mathcal{I}$（第 1 章定义，对象改为 $\vect{z}$）。
\item 若 $\Delta\mathcal{I}\le 0$ 且 $\alpha_m$ 的 $5\%$ 分位已高于安全线，不探针，进入停机或保温判断。
\item 否则对 $\cA=\{P0,\ldots,P4\}$ 用当前集合向前滚 $\tau(a)$，计算 $G$ 与 $C$。滚动时执行器饱和、传感器饱和按真实约束写入。
\item 选取 $a^\star$。若最优仍是 P0，保持加工轨迹，并提高对用户的不确定度提示。
\item 执行探针期间冻结慢速菜单逻辑，快安全环继续工作。
\item 探针结束后用新观测更新 $\vect{\theta}$ 与 $\vect{z}$。若 $\det\vect{\Sigma}$ 未下降到预测的 $70\%$ 以内，标记模型失配，切换到更保守的机会约束。
\end{enumerate}

第 6 步是防止“探针 theater”：动作做了，信息没有进来，就不允许控制器宣称已经认知。

\section{对偶目标}

探针库解决的是家电执行器不能任意波形。决策层仍把信息与品质放在同一次优化里，避免探针状态机和加工状态机抢执行器。

\begin{equation}
\label{eq:dual}
\begin{aligned}
\max_{\vect{u}_{k:k+H}}\;
& \E\big[Q(\vect{z}_{k+H})\big]
+\lambda_I\, G(\vect{u}_{k:k+H})
-\lambda_E E(\vect{u})\\
\mathrm{s.t.}\quad
& \Delta\alpha_b\le\eta,\quad
\Delta m\le m_{\mathrm{budget}},\quad
\vect{u}\in\cU.
\end{aligned}
\end{equation}
$Q$ 的定义在第 \ref{ch:safe} 章。$\lambda_I$ 随 $\Delta\mathcal{I}$ 增大。赤字低于零且安全分位已通过时，$\lambda_I=0$，最优解退回纯品质控制。实现时先在探针库上比较 $G-\lambda C$，再把胜出的探针作为 \eqref{eq:mpc} 控制时域中的一段固定输入。连续最优波形只作为离线对照，用来检查探针库损失了多少信息。

\section{短脉冲几乎不推动冷点}

表面节点时间常数 $\tau_1=R_{\mathrm{ext}}C_1$，内部时间常数 $\tau_2=R_{\mathrm{int}}C_2$。在 $\tau\ll\tau_2$ 的窗口内，冷点温度的增量是 $O(\tau/\tau_2)$。由 \eqref{eq:F}，冷点进度增量同样是这个量级。损伤预算因此主要约束表面褐变和表面含水，而不是冷点安全。厚度增大使 $\tau_2$ 增大，窗口更宽，但表面响应对 $R_{\mathrm{int}}$ 的灵敏度下降。信息增益存在由能量、$\eta$ 和 $\tau_{\max}$ 共同限制的上界。

仿真在厚度与初始含水平面上计算该上界。上界低于停机所需的熵减时，界面输出进度区间，不输出一个剩余时间点。显示宽度与后验预测区间一致。

熵与 Fisher 的离散推导、探针灵敏度矩阵和该等高线的计算步骤见附录。
